% Proposed author-reviewed replacement for the short Research Positioning section.
% Insert corresponding bibliography items from REFERENCES_TO_ADD.tex.
\section{Related work and contribution boundaries}\label{sec:related-work-expanded}

\paragraph{Zero-error coding with side information.}
Witsenhausen's zero-error side-information framework relates minimum
message alphabets to graph colorings~\cite{witsenhausen1976}. In the static
case of our model, two binary words can share an encoding label precisely
when their Hamming distance exceeds the erasure budget; the resulting
chromatic number of a power of the binary cube is the appropriate baseline.
Exact finite colorings, including Kokkala and \"Osterg\r{a}rd's solution
for the square of the eight-cube~\cite{kokkala2017}, demonstrate that
these static invariants need not admit simple uniform formulas. More
recent graph-entropy formulations of zero-error computation with
side information further illustrate the breadth of the static
characteristic-graph approach~\cite{charpenay2024}. Our additional
constraint is a consistent transition mapping for every specified blind
coordinate assignment. Thus an initial color class cannot be optimized
without regard to the labels available after a write.

\paragraph{Causal coding and rewritable storage.}
Witsenhausen also investigated the structure of finite-horizon real-time
source coders under causal information constraints~\cite{witsenhausen1979}.
Ahlswede and Zhang's write-efficient memory model studies data rewriting
through a transition cost and its associated rate--cost tradeoff
\cite{ahlswedezhang1989}. These antecedents motivate the distinction
between constraints on the encoder's information, the number of encoded
states, and the cost of changing a physical representation. Our model
charges the cardinality of a persistent state alphabet, makes the
previously stored source coordinate unavailable to the online transition,
and supplies the current unerased coordinates to the decoder. It does not
optimize the Hamming distance between successive physical codewords.

\paragraph{Locally updatable codes and oblivious updates.}
Chandran, Kanukurthi and Ostrovsky introduced locally updatable and
locally decodable codes that balance limited codeword modifications and
recovery locality~\cite{chandran2014}. Nakkiran, Shah and Rashmi derived
communication bounds and constructions for oblivious updates of stale
storage nodes~\cite{nakkiran2014}. Both concern changes to encoded
storage, but the latter counts refresh communication and the former
counts locality/rewrite operations; neither has the identical one-write,
zero-error, phase-aware state-alphabet objective studied here.

\paragraph{Specific contributions and their scope.}
Under the explicit assumption of one blind coordinate assignment, a
public phase index, and zero-error recovery for all erasure patterns,
we derive a high-erasure matching normal form and bound the multiplicity
of forced successor edges by initial hypercube adjacencies. Combining
these tools with explicit cubical constructions gives the exact
state-alphabet plateau at the two adjacent erasure budgets $q=n-3$
and $q=n-2$. The exceptional seven-coordinate lower bound uses an
exhaustive finite reduction and separately checkable propositional
refutations. The statement concerns exact \emph{label counts}; the rounded
optimal bit budget remains $n$ in the stated regime. The analysis does
not establish a complexity classification for arbitrary update horizons
or a general coding theorem for all dynamic-storage models.